% Contenuto della «Bussola in 7 domande».
% Usato sia nell'appendice del libro (appendix/bussola.tex)
% sia nella versione stampabile autonoma (misc/bussola_stampabile.tex).
\begin{enumerate}[label={\sffamily\bfseries\large\arabic*}, leftmargin=2.2em, itemsep=0.9em]
	\item \textbf{Prove contrarie.} \textit{Che cosa mi farebbe cambiare idea?} Se non riesco a immaginare nessuna prova capace di smentirmi, non sto ragionando: sto difendendo.

	\item \textbf{Probabilità.} \textit{Quanto è frequente, in generale, quello che sto valutando?} Prima i numeri, poi la storia: il caso che mi colpisce somiglia allo stereotipo, ma quante volte accade davvero?

	\item \textbf{Fonte.} \textit{Da dove viene questa informazione, e come lo saprebbe?} Uno studio, un'esperienza personale, un'opinione, un testo generato da una macchina? Una risposta fluida non è una risposta verificata.

	\item \textbf{Incentivo.} \textit{Chi guadagna se ci credo, o se la condivido?} Vale per chi mi parla, per la piattaforma che me la mostra e per me stesso: a volte l'incentivo più forte è il desiderio di avere ragione.

	\item \textbf{Costo dell'errore.} \textit{Che cosa succede se sbaglio, e quanto è reversibile?} Un errore su un film costa poco; uno sulla salute, sul denaro o su una persona può costare molto. Più alta la posta, più lenta la decisione.

	\item \textbf{Tempo.} \textit{Devo decidere adesso?} Stanchezza, fretta e indignazione sono le condizioni in cui le scorciatoie vincono. Se posso, rimando di un'ora, di una notte.

	\item \textbf{Confronto.} \textit{Chi, tra le persone che rispetto, la vede diversamente?} Chiedo il suo argomento più forte, non la sua versione più facile da smontare. È la correzione che da solo non vedrei.
\end{enumerate}
